\documentclass[11pt]{article}

% --- Packages ---
\usepackage[utf8]{inputenc}
\usepackage[T1]{fontenc}
\usepackage{lmodern}
\usepackage[margin=1in]{geometry}
\usepackage{microtype}

\usepackage{graphicx}
\usepackage{xcolor}
\usepackage{amsmath, amssymb}
\usepackage{siunitx}
\usepackage{booktabs}
\usepackage{caption}
\usepackage[numbers,sort&compress]{natbib}
\usepackage{hyperref}
\hypersetup{colorlinks=true, linkcolor=black, citecolor=black, urlcolor=blue}

\usepackage{authblk}

\graphicspath{{figures/pdf/}}

% --- Shared color palette ---
% Edit these once; use everywhere via \color{cleoBlue} etc.
\definecolor{cleoBlue}{HTML}{2E5C8A}
\definecolor{cleoOrange}{HTML}{D97A29}
\definecolor{cleoGreen}{HTML}{4A8C5A}
\definecolor{cleoRed}{HTML}{B23A3A}
\definecolor{cleoGray}{HTML}{6B7280}

% --- Convenience macros ---
\newcommand{\cleo}{\textsc{cleo}}
\newcommand{\todo}[1]{\textcolor{cleoRed}{\textbf{[TODO: #1]}}}


% ---------------------------------------------------------------------------
% The lab notebook.
%
% This is deliberately a SEPARATE document from main.tex. The two have
% different readers and different standards of evidence:
%
%   main.tex      the paper. Only claims that survived replication, written as
%                 prose for a referee. Nothing enters it that is not settled.
%   notebook.tex  the run log. Every experiment E1--E21, including the ones
%                 that failed, the ones that were retracted, and the reasoning
%                 that turned out to be wrong. Pre-registrations stay next to
%                 the results that falsified them.
%
% Keeping the record is the point of the split, not tidiness. E18's frontier
% was retracted by E19 using E18's own pre-registered criterion; that only
% works because the criterion was written down before the data landed. A paper
% cannot carry that and stay readable, and deleting it would destroy the
% evidence that the retraction was principled rather than post-hoc.
%
% Build:  make notebook   ->  notebook.pdf
% ---------------------------------------------------------------------------

\title{\cleo{}: experimental record}
\author{}
\date{Lab notebook --- compiled \today}

\begin{document}

\maketitle
\thispagestyle{empty}

\begin{abstract}
\noindent
The complete experimental record behind \cleo{}: arm definitions, selection
algorithms, the E1--E21 slate with pre-registrations and results, the algorithm
audit, and the scoring protocol. Retracted claims are kept in place with the
reasoning that produced them, rather than removed. This document is the
provenance for the numbers quoted in the paper; it is not itself a paper.
\end{abstract}

\clearpage
\setcounter{tocdepth}{2}
\tableofcontents
\clearpage

\section{Experimental programme and selection algorithms}
\label{sec:experiments}

\plan{This section exists so that every arm, rule and planned experiment is
written down in one place with a fixed name, and so that results reported
elsewhere in the paper can be traced to a defined procedure. Much of it will
move to Materials and Methods at submission; keeping it in the main build for
now makes disagreements visible early.}

\subsection*{Notation}

A design is a sequence $s \in \mathcal{A}^L$ over the 20 canonical amino acids.
Distance is Hamming,
\[
  d(s, t) \;=\; \bigl|\{\, i : s_i \neq t_i \,\}\bigr|,
\]
which is the metric behind every diversity number we report. One-hot encoding
$\phi(s) \in \{0,1\}^{20L}$ satisfies
$\lVert \phi(s) - \phi(t) \rVert_2 = \sqrt{2\, d(s,t)}$, so the PCA projections
in the figures are a linear embedding of exactly this distance rather than a
second, competing notion of similarity. Where an ESM embedding is used instead,
$d$ is replaced by Euclidean distance between mean-pooled representations; this
is reported as a diagnostic and, optionally, as a secondary library metric.

A design \emph{passes} when its motif all-atom RMSD is below
\SI{1.5}{\angstrom} and no backbone atom clashes with the ligand. $\mathcal{P}$
denotes the passing subset of a pool $\mathcal{S}$. The
\emph{anchor} $a$ is the position-wise consensus of low-temperature LigandMPNN
samples --- the single sequence best representing the mode the base model's
likelihood concentrates on.

\subsection*{Selection rules}

\plan{All implemented in \texttt{src/cleo/design/utils/selection.py}, all
sharing one distance backend so the comparison is between rules rather than
between implementations. Benchmarked by
\texttt{paper/figures/ame\_selection\_bench.py}.}

Given a pool $\mathcal{S}$ of $M$ sampled sequences and a budget of $k \ll M$
folds, each rule returns the subset to fold.

\begin{description}
  \item[\texttt{random}] Uniform without replacement. \textbf{This is the
        control, and it is the bar.} Diversity is trivially won and worthless
        alone: a rule that doubles spread while halving passing designs has made
        the library worse.
  \item[\texttt{maxmin}] Greedy farthest-point (k-center) traversal. Having
        selected $\mathcal{T}$, repeatedly add
        $\arg\max_{s \in \mathcal{S} \setminus \mathcal{T}}
         \min_{t \in \mathcal{T}} d(s,t)$.
        The inner \emph{minimum} is what distinguishes this from max-sum
        selection: it prevents accepting a tight cluster that happens to sit far
        from everything already chosen.
  \item[\texttt{maxmin\_anchor}] As above, scoring
        $w_{\mathrm{self}} \min_{t \in \mathcal{T}} d(s,t) + w_{\mathrm{anchor}}\, d(s,a)$.
        Selects designs novel relative to the incumbent low-temperature
        solution, not merely relative to the rest of the library.
  \item[\texttt{anchor\_far}] Rank by $d(s,a)$ descending. The pure
        ``maximally unlike vanilla MPNN'' rule.
  \item[\texttt{anchor\_band}] Restrict to the shell
        $\{\, s : q_{\mathrm{lo}} \le d(s,a) \le q_{\mathrm{hi}} \,\}$ for
        quantiles of the pool's anchor-distance distribution, then run
        \texttt{maxmin} inside it. Designed as the correction to
        \texttt{maxmin}: bounded above as well as below, so the traversal cannot
        walk out to sequences that are distant because they are degenerate.
  \item[\texttt{cluster\_rep}] Average-linkage clustering cut to $k$ clusters,
        one representative sampled per cluster. Draws from the pool's density
        rather than its edges.
\end{description}

\subsection*{Result: diversity-first acquisition does not work}

\plan{Reported here rather than buried, because it was a starred hypothesis and
because the negative result is informative. Retrospective evaluation on the
M0097 trajectory pool: every sampled sequence already carries a pass label, so
any rule can be scored without folding anything.}

Passing designs recovered, relative to \texttt{random}, at a budget of 400 folds:

\begin{center}
\begin{tabular}{@{}lc@{}}
\toprule
\textbf{Rule} & \textbf{Passing designs / random} \\
\midrule
as-sampled          & 2.04 \\
\texttt{random}     & 1.00 \\
\texttt{cluster\_rep}   & 0.40 \\
\texttt{anchor\_band}   & 0.36 \\
\texttt{maxmin}         & 0.36 \\
\texttt{maxmin\_anchor} & 0.24 \\
\texttt{anchor\_far}    & 0.12 \\
\bottomrule
\end{tabular}
\end{center}

Every diversity-directed rule loses to uniform random selection, by factors of
\numrange{2.5}{8}. The mechanism is visible in the spread: no rule moves mean
pairwise distance beyond \numrange{135}{144} residues against random's 135. The
pool is already near-saturated on diversity, so a selection rule has almost
nothing to gain and a great deal to lose --- any systematic deviation from
random correlates with selecting sequences that fail. Banding helps relative to
raw farthest-point traversal ($0.36$ vs.\ $0.07$ at $k = 200$) but does not
approach the control.

\plan{Two honest caveats. The as-sampled row beats random because the pool is
time-ordered and later policies are better --- that is a policy-quality effect,
not a selection effect, and it must not be reported as one. And absolute counts
are small (33 passing in a 1200-design pool), so individual cells are noisy;
what carries the conclusion is that the direction is consistent across six
rules and three budgets.}

\textbf{Consequence for the paper.} Selection is not where the gain is. The
defensible use of these rules is post-hoc: choosing a physically orderable
subset from designs \emph{already known to pass}, where the labels exist and
only the diversity trade-off is live. Deciding what to fold should be left to
the policy.

\subsection*{Experiment slate}

\plan{Status as of the pilot. \statusdone{} = measured, \statusscaffold{} =
partial, \statusblocked{} = not started. Full protocol detail lives in
\texttt{paper/experiments\_plan.md}. The \textbf{Result} column is filled in
as each run lands, so the table doubles as the run log.

E5 (held-out oracle) and E6 (policy transfer) are \emph{deliberately deferred},
not dropped. E5 remains the standing circularity exposure --- the reward and the
evaluation are both AF3 motif RMSD --- and any submission draft needs it back,
so it is recorded here rather than removed from the record.}

\begin{center}
\begin{tabular}{@{}llp{0.40\textwidth}p{0.22\textwidth}@{}}
\toprule
\textbf{ID} & \textbf{Status} & \textbf{Experiment} & \textbf{Result} \\
\midrule
E1  & \statusscaffold & Matched-fold head-to-head &
      M0097 per 8 folds: $T{=}0.1$ 5.0 passing, \cleo{} 3.1 --- but see E10,
      the comparison must be made at matched diversity \\
E2  & \statusscaffold & Diversity yield curves &
      \cleo{} 3200 folds $\rightarrow$ 97 passing $\rightarrow$ 97 clusters,
      still climbing \\
E3  & \statusblocked  & Rescue of the three near-miss backbones
      (\SIrange{1.61}{1.69}{\angstrom} best-of-40, 0/40 published) & \todoresult \\
E4  & \statusblocked  & Compute efficiency; cross-cutting over E1--E3 & \todoresult \\
E7  & \statusscaffold & Effective library size &
      Among passing designs, cluster count equals design count \\
E8  & ---             & Mutation-diversity reward &
      \textbf{Retired}, superseded by trajectory-as-library \\
E9  & \statusdone     & Selection rules, retrospective &
      \textbf{Negative:} every rule loses to random, \numrange{2.5}{8}$\times$ \\
\addlinespace
E10 & \statusdone     & \textbf{Temperature sweep} $T \in \{0.1, \ldots, 1.0\}$;
      builds the incumbent frontier &
      \textbf{Frontier shift confirmed:} at matched diversity \cleo{} passes
      $1.75\times$ the baseline; at matched pass rate it is $1.46\times$ more
      diverse \\
E11 & \statusblocked  & Reward shaping: rank normalisation & \todoresult \\
E12 & \statusblocked  & Early stopping at ${\sim}75$ steps & \todoresult \\
E13 & \statusblocked  & Consensus anchoring of the search & \todoresult \\
E14 & \statusdone     & ESM embedding view of the passing libraries &
      Advantage holds but shrinks: passing-set spread \cleo{}/$T{=}0.1$ is
      $1.31\times$ in ESM space vs.\ $2.57\times$ in Hamming \\
E15 & \statusblocked  & Centering-weight sweep, $W$ in both directions & \todoresult \\
E16 & \statusblocked  & In-training selection: no-oversample /
      oversample+random / oversample+diverse & \todoresult \\
\bottomrule
\end{tabular}
\end{center}

\todo{experiments --- run E10; it gates the central claim}

\subsection*{E15 --- Centering weight: one knob from centered to diverse}

\plan{No new reward machinery. \texttt{compute\_dist\_to\_ref\_seqs\_from\_df}
already emits Hamming distance to a reference set, and the aggregator already
supports per-metric \texttt{mode}, \texttt{weight} and rank normalisation, so
the whole experiment is a config sweep. Configs generated by
\texttt{experiments/ame/make\_centering\_sweep.py}.}

The reference $a$ is the consensus of low-temperature LigandMPNN samples. Adding

\begin{center}
\texttt{- metric: dist\_to\_ref\_min \quad mode: \{min|max\} \quad weight: $W$ \quad normalize: rank}
\end{center}

to the aggregation gives a single signed axis. Because reward is a weighted
average, $\smash{R = \sum_i w_i r_i / \sum_i w_i}$, the weight is directly
interpretable: $W = 0$ is pure motif RMSD, $W = 1$ gives centering equal
influence to it. \texttt{mode:\ min} pulls the policy toward the incumbent
mode; \texttt{mode:\ max} pushes it away.

Sweep $W \in \{0, 0.25, 0.5, 1.0\}$ in both directions. Prediction: mean
distance to $a$ falls monotonically under \texttt{min} and rises under
\texttt{max}, and the PCA panels show the passing cloud sliding toward or away
from the low-temperature knot. Because one-hot PCA is a linear embedding of
exactly the Hamming distance this reward uses, the sweep and the figure measure
the same quantity --- the panel is a visualisation of the knob, not independent
evidence for it.

\plan{The knob will move designs; that is not in question and is not the
result. What the experiment has to establish is whether \emph{pass rate
survives} at either end. Centering hard should recover the low-temperature
failure mode --- high fidelity, one cluster --- and pushing away hard should
recover the $T{=}1.0$ failure mode --- broad and dead. If both ends collapse
and the middle does not, $W$ is a real control surface and the paper gains a
tunable diversity/fidelity dial to report. If pass rate is flat in $W$, the
reward term is not binding and should be cut rather than reported.

This also supplies the honest version of the frontier: E10 sweeps the
\emph{baseline's} knob (temperature) and E15 sweeps \emph{ours}. Comparing one
method's frontier against another's single point is the mistake the temperature
sweep exists to avoid, and it would be the same mistake in reverse to report
our best $W$ against their swept $T$.}
\todo{experiments --- run E15 on M0097 once E11 rank-normalisation is validated}

\subsection*{E14 result --- the embedding view qualifies the Hamming view}

\plan{Mean-pooled ESM-2 650M, one vector per design, computed for every arm.}

Within a set, ESM distance and Hamming distance are close to \emph{uncorrelated}
($r$ between $-0.17$ and $0.56$ across arms), so the embedding is genuinely
measuring something the Hamming numbers do not. That makes the agreement in
ordering meaningful rather than automatic: on M0097's passing designs,
\cleo{} is more spread than $T{=}0.1$ under both metrics.

But the size of the advantage does not survive the change of metric. The ratio
of \cleo{}'s passing-set spread to the low-temperature arm's is
$\mathbf{2.57\times}$ in Hamming and only $\mathbf{1.31\times}$ in ESM space.
A substantial share of the extra sequence distance \cleo{} buys is
biochemically conservative --- positions that differ in identity but not much
in what ESM considers the residue to be doing.

\plan{This should be stated in the paper, not left for a referee to find.
Reporting ``2.6$\times$ more diverse'' with a metric chosen because it is the
larger number is exactly the move the temperature sweep exists to prevent us
from making about temperature. The defensible claim is that the advantage is
robust to the choice of metric while its magnitude is not, and that the
library-relevant quantity --- distinct passing clusters per fold, where
\cleo{} wins $3.1$ to $1.0$ --- does not depend on either distance at all,
since those designs are separate clusters at any threshold tested.

Caveat on $n$: passing subsets are 5 and 15 designs, so the correlation
figures in particular are noisy. The ratio comparison is the robust part.}

\subsection*{E10 result --- \cleo{} sits above the temperature frontier}

\plan{The experiment the section was gated on. LigandMPNN sampled at six
temperatures, 24 sequences per backbone, folded best-of-5, and scored by the
same code path as \cleo{}. M0097, at the benchmark's 8-fold budget, averaged
over 300 random draws of 8.}

\begin{center}
\begin{tabular}{@{}lcccc@{}}
\toprule
\textbf{Arm} & \textbf{Mean pairwise} & \textbf{Variable} &
\textbf{Passing} & \\
             & \textbf{Hamming}       & \textbf{positions} & \textbf{/ 8 folds} & \\
\midrule
LigandMPNN $T{=}0.1$ & 39.1  & 81    & 5.00 & \\
LigandMPNN $T{=}0.2$ & 48.7  & 102   & 4.81 & \\
LigandMPNN $T{=}0.3$ & 58.7  & 116   & 4.42 & \\
LigandMPNN $T{=}0.5$ & 78.2  & 135   & 2.23 & \\
LigandMPNN $T{=}0.7$ & 104.8 & 159   & 1.72 & \\
LigandMPNN $T{=}1.0$ & 123.8 & 170   & 0.00 & \\
\midrule
\textbf{\cleo{}-GRPO} & \textbf{103.3} & \textbf{163} & \textbf{3.06} & \\
\bottomrule
\end{tabular}
\end{center}

Temperature traces a clean monotone trade-off: diversity rises from 39 to 124
mean pairwise substitutions while the pass rate falls from 5.0 to 0.0 per eight
folds. That curve is the incumbent frontier, and it is the thing every claim in
this section has to beat.

\cleo{} sits above it. Interpolating the baseline curve:

\begin{itemize}\setlength{\itemsep}{0pt}
  \item at \cleo{}'s diversity (Hamming 103.3), the baseline passes
        $1.75$ per 8; \cleo{} passes $\mathbf{3.06}$ --- $\mathbf{1.75\times}$.
  \item at \cleo{}'s pass rate ($3.06$ per 8), the baseline reaches Hamming
        $70.8$; \cleo{} reaches $\mathbf{103.3}$ --- $\mathbf{1.46\times}$.
\end{itemize}

The two dangerous points named in advance, $T{=}0.2$ and $T{=}0.3$, do not
threaten the claim: they buy very little diversity (48.7 and 58.7) for the pass
rate they retain, landing far below \cleo{}'s coverage. The nearest baseline
arm by diversity is $T{=}0.7$, which matches \cleo{} on both diversity measures
(104.8 vs.\ 103.3 Hamming; 159 vs.\ 163 variable positions) and passes
\emph{half} as often.

\plan{\textbf{Correction to carry through the rest of the paper.} Earlier
drafts reported ``$T{=}0.1$ yields 1.0 cluster per 8 folds against \cleo{}'s
3.1''. That used a \SI{70}{\percent}-identity clustering threshold; at
\SI{90}{\percent} identity every arm's passing designs are separate clusters
and the cluster count simply equals the passing count. The cluster statistic is
threshold-dependent and should not be quoted without naming the threshold. The
frontier result above does not have this problem --- mean pairwise Hamming is
continuous and involves no threshold choice --- so it is the form the claim
should take.}


\clearpage
\addcontentsline{toc}{section}{References}
\bibliographystyle{unsrt}
\bibliography{refs}

\end{document}
